\cleardoublepage

\chapter{Podsumowanie}
Cel niniejszej pracy dyplomowej został spełniony. Dzięki autorskiemu oprogramowaniu do automatyzacji testów, zbadano 30 konfiguracji systemów RAG wykorzystujących pięć najnowszych popularnych modeli językowych firm takich jak: OpenAI, DeepSeek, Google, Anthropic oraz xAI. Ocena wyników opierała się na następujących kryteriach: trafność odpowiedzi, koszt i czas wygenerowania odpowiedzi oraz przepustowość tokenów. Ponadto, poprzez intuicyjny graficzny interfejs użytkownika, zaimplementowano dwa najwydajniejsze systemy RAG: jeden przeznaczony do danych nieustrukturyzowanych oraz drugi do danych ustrukturyzowanych.

Wyniki badania okazały się być zaskakująco dobre, ale nieporównywalne z rezultatami trzech omawianych artykułów naukowych. Uważa się, że istnieje jedna przyczyna takich różnic. Autor w ramach testów odpytał systemy RAG za pomocą łącznie szesnastu autorskich pytań. Natomiast, w publikacjach twórcy użyli ich 360, 1000 czy nawet 4780.

Jeżeli prace w ramach tego tematu byłyby kontynuowane, proponuje się poszerzenie bazy pytań do co najmniej pięciuset, żeby dorównać podejściu stosowanemu w publikacjach. Niesie to za sobą potrzebę zastosowania automatyzacji w następujących obszarach: tworzenie pytań, przekazywanie ich do dedykowanego oprogramowania oraz analiza wyników. Dla pierwszego i trzeciego podejścia niezbędny będzie odrębny LLM i oprogramowanie weryfikujące poprawność odpowiedzi, tak jak to przedstawiono w publikacji \cite{art1_2025} w ramach benchmarku MEBench.
